% =============================================================
% Capítulo 4 — Metodologia
% =============================================================

\chapter{Metodologia}\label{chap:metodologia}

Este capítulo apresenta o percurso metodológico adotado no trabalho, organizado em cinco
etapas: a seleção do modelo de referência, a implementação das estratégias de poda sem
otimização, a implementação da estratégia de poda com otimização, o \textit{fine-tuning} de
recuperação e a definição das métricas de avaliação. Para cada etapa, estabelecem-se os
critérios que orientaram as decisões tomadas. As decisões concretas decorrentes desses
critérios --- incluindo o modelo selecionado, o \textit{dataset} empregado e as
configurações de poda --- são descritas no Capítulo~\ref{chap:desenvolvimento}.

% -------------------------------------------------------
\section{Seleção e Implementação do Modelo de Referência}\label{sec:metodologia_modelo}
% -------------------------------------------------------



O modelo de referência (\textit{baseline}) corresponde à versão original do modelo de linguagem
de grande escala (\textit{Large Language Model} --- LLM) antes de qualquer processo de poda.
Ele serve como padrão para mensurar a perda de desempenho introduzida pela compressão: todas as
métricas dos modelos podados são comparadas às métricas desse modelo original.

Para a seleção do modelo de referência, foram adotados os seguintes critérios:

\begin{itemize}
    \item ser um LLM representativo, amplamente utilizado na literatura de compressão de modelos;
    \item dispor de pesos pré-treinados publicamente disponíveis;
    \item apresentar tamanho compatível com as restrições de \textit{hardware} disponível;
    \item possuir ampla documentação em trabalhos de poda neural;
    \item \modificado{contar com implementação acessível em \textit{frameworks} consolidados,
          como PyTorch.}
\end{itemize}

% AJUSTE: reunião 13/05 — Hugging Face como plataforma, não biblioteca
\adicionadomarcelo{A implementação do modelo de referência utiliza a biblioteca PyTorch \cite{Paszke2019}.
Como repositório de modelos pré-treinados, é empregada a plataforma Hugging Face, um
ambiente colaborativo que disponibiliza ampla variedade de modelos de linguagem de grande
escala publicamente acessíveis --- como, por exemplo, variantes do GPT-2.
O modelo selecionado a partir dos critérios estabelecidos --- incluindo a versão específica
obtida via Hugging Face --- é apresentado e detalhado no
Capítulo~\ref{chap:desenvolvimento}.}

% -------------------------------------------------------
\section{Implementação das Estratégias de Poda sem Otimização}\label{sec:metodologia_poda_sem}
% -------------------------------------------------------


% AJUSTE: reunião 13/05 — duas estratégias (magnitude + estruturada), "avaliadas" em vez de "consideradas"
\adicionadomarcelo{As estratégias de poda sem otimização têm por objetivo reduzir o número de parâmetros de um
modelo pré-treinado, eliminando pesos ou estruturas avaliadas como menos relevantes para o
desempenho preditivo. Serão avaliadas duas estratégias de poda clássicas, escolhidas por
representarem os dois extremos do espectro de granularidade da poda neural:

\begin{itemize}
    \item poda baseada em magnitude dos pesos \cite{Han2015}, que remove parâmetros
          individuais cujos valores absolutos são inferiores a um limiar; representa a
          granularidade mais fina possível e gera modelos esparsos sem alteração da
          estrutura física do modelo;
    \item poda estruturada \cite{Li2017}, que elimina unidades arquiteturais inteiras
          --- filtros convolucionais, no trabalho original; em uma arquitetura
          \textit{Transformer}, cabeças de atenção ou neurônios da camada intermediária do
          MLP --- reduzindo fisicamente as dimensões dos tensores e gerando modelos
          computacionalmente mais enxutos.
\end{itemize}}

% AJUSTE: reunião 13/05 — reforçar que estratégias são consolidadas na literatura
\adicionadomarcelo{Ambas as estratégias são amplamente consolidadas na literatura de compressão de modelos
neurais, com mais de uma década de aplicação e extensa validação experimental. Sua inclusão
neste trabalho não constitui contribuição metodológica original; pelo contrário, sua
aplicação fornece uma base de comparação estabelecida, sobre a qual é avaliada a
estratégia otimizada proposta na
Seção~\ref{sec:metodologia_poda_com}.}

Os critérios que orientaram a escolha das estratégias a implementar foram:

\begin{itemize}
    \item representatividade na literatura de compressão de modelos;
    \item complementaridade entre granularidades distintas de remoção;
    \item viabilidade de implementação com as ferramentas disponíveis.
\end{itemize}


Em ambas as estratégias, são aplicadas funções de avaliação de importância de parâmetros,
que orientam o processo de remoção: a poda por magnitude utiliza como critério o valor
absoluto de cada peso individual, enquanto a poda estruturada utiliza a norma $L_1$ do
conjunto de pesos associado a cada unidade arquitetural removível.



As estratégias selecionadas com base nesses critérios, bem como os detalhes de sua
implementação, são apresentados no Capítulo~\ref{chap:desenvolvimento}.

% -------------------------------------------------------
% AJUSTE: reescrita da Seção 4.3 — a estratégia otimizada passa a ser declarada
% (sensibilidade por camada); taxas dinâmicas, momento de aplicação e Knowledge
% Distillation descem para alternativas descartadas, como no Cap. 5.
\section{Implementação da Estratégia de Poda com Otimização}\label{sec:metodologia_poda_com}
% -------------------------------------------------------

Esta é a etapa central do trabalho. O objetivo é investigar se a combinação de uma
estratégia clássica de poda com uma técnica de otimização resulta em modelos mais
eficientes do que a poda aplicada de forma direta, sem qualquer refinamento adicional.

A técnica de otimização adotada é a \textbf{análise de sensibilidade por camada}, que
orienta uma alocação não uniforme do orçamento de poda entre as camadas do modelo. A
hipótese que ela investiga é a de que as camadas não são igualmente importantes: se
algumas toleram a remoção de parâmetros muito mais do que outras, distribuir o orçamento
de poda em proporção a essa tolerância deve produzir, para um mesmo grau de compressão,
perda menor de desempenho preditivo do que as alocações ingênuas adotadas pelas
estratégias da Seção~\ref{sec:metodologia_poda_sem}.

O procedimento organiza-se em duas fases. Na primeira, levanta-se um perfil de
sensibilidade do modelo: poda-se uma camada de cada vez, a uma taxa fixa e mantendo todas
as demais intactas, e registra-se a degradação resultante na métrica de desempenho
preditivo. Na segunda, esse perfil é convertido em uma taxa de poda por camada, atribuída
em proporção à robustez medida e escalada de modo que o orçamento global pretendido seja
respeitado. Como cada estratégia de poda remove unidades de natureza distinta, o perfil é
levantado separadamente para cada uma delas. A formulação da regra de alocação, a
taxa-sonda empregada no perfilamento e os valores adotados para os demais parâmetros são
apresentados no Capítulo~\ref{chap:desenvolvimento}.

Os critérios que orientaram a escolha dessa técnica entre as alternativas cogitadas foram
sua aderência direta à hipótese do trabalho e o fato de não alterar a natureza da poda:
ela redistribui o orçamento entre as camadas sem modificar o critério de importância nem
introduzir um segundo modelo no experimento. Isso mantém a comparação com as estratégias
sem otimização restrita a uma única variável.

\subsection*{Alternativas consideradas e descartadas}

Outras técnicas de otimização foram cogitadas e não integram o escopo experimental, por
restrição de tempo. Registram-se a seguir, por delimitarem o alcance das conclusões e por
constituírem desdobramentos naturais para trabalhos futuros:

\begin{itemize}
    \item \textbf{comparação entre taxas de poda fixas e dinâmicas}, em que o orçamento de
          poda varia ao longo do processo em vez de ser fixado de antemão;
    \item \textbf{variação do momento de aplicação da poda} --- antes, durante ou após o
          treinamento. Este trabalho aplica a poda sempre sobre o modelo pré-treinado,
          reservando o retreinamento para a etapa de recuperação descrita na
          Seção~\ref{sec:metodologia_finetuning};
    \item \textbf{\textit{Knowledge Distillation}}, na qual um modelo menor
          (\textit{student}) é treinado para reproduzir o comportamento de um modelo maior
          (\textit{teacher}); trata-se de técnica de compressão de natureza distinta da
          poda, cuja avaliação exigiria protocolo experimental próprio;
    \item \textbf{quantização}, que reduz a precisão numérica da representação dos pesos
          em vez de eliminá-los.
\end{itemize}

A comparação entre os modelos obtidos com e sem otimização visa identificar em que medida
a alocação guiada por sensibilidade melhora o balanço entre desempenho preditivo e
eficiência computacional em relação às alocações ingênuas.

% -------------------------------------------------------
% AJUSTE: seção nova — o fine-tuning de recuperação era, até aqui, apenas um
% sub-item da lista de estratégias de otimização, embora seja uma etapa própria
% no Cap. 5 e uma seção inteira do Cap. 6.
\section{\textit{Fine-tuning} de Recuperação}\label{sec:metodologia_finetuning}
% -------------------------------------------------------

As etapas anteriores operam em regime \textit{one-shot}: o modelo é podado e medido
imediatamente, sem retreinamento. Esse regime é deliberado, pois isola o dano causado pela
poda do efeito de qualquer recuperação posterior, mas não corresponde ao procedimento da
literatura --- tanto \citeonline{Han2015} quanto \citeonline{Li2017} retreinam o modelo
depois de podá-lo. Esta etapa fecha essa lacuna, medindo quanto da degradação é recuperável
sob um orçamento de retreinamento limitado.

Dois cuidados metodológicos condicionam a validade dessa medição.

O primeiro é a necessidade de uma \textbf{linha de controle densa}. O modelo de referência
não foi pré-treinado no corpus de avaliação, de modo que retreiná-lo nesse corpus já
melhora o desempenho preditivo por simples adaptação de domínio, sem poda alguma
envolvida. Sem essa linha de controle, seria impossível distinguir a recuperação do dano da
poda da mera adaptação ao corpus. O modelo denso é, portanto, submetido ao mesmo orçamento
de retreinamento e avaliado junto às demais configurações.

O segundo é a \textbf{preservação da esparsidade}. Na poda não estruturada os pesos são
apenas anulados, e o otimizador os restauraria ao longo do treinamento, desfazendo a
compressão que se pretende avaliar. Como em \citeonline{Han2015}, a esparsidade é mantida
por uma máscara binária fixa, capturada imediatamente após a poda e reaplicada a cada passo
do otimizador. Na poda estruturada esse cuidado é dispensável, pois as unidades removidas
deixaram de existir.

Registre-se, por fim, que esta é a única etapa estocástica do experimento: as varreduras
\textit{one-shot} são determinísticas, e a semente aleatória passa a afetar o resultado
apenas aqui. O custo energético do retreinamento é medido separadamente do custo de
inferência, por se tratar de custo pago uma única vez, que precisa ser amortizado contra a
economia obtida a cada inferência.

O protocolo concreto --- otimizador, orçamento de passos, hiperparâmetros e configurações
submetidas ao retreinamento --- é apresentado no Capítulo~\ref{chap:desenvolvimento}.

% -------------------------------------------------------
% AJUSTE: reunião 13/05 — reescrita completa da Seção 3.4 (métricas explicadas em detalhe)
\section{Métricas de Avaliação}\label{sec:metodologia_metricas}
% -------------------------------------------------------

\adicionadomarcelo{As métricas de avaliação são organizadas em três dimensões: desempenho preditivo, custo
computacional e impacto ambiental. Para cada métrica, descreve-se o que ela representa e
como é obtida operacionalmente.

\subsection*{Desempenho preditivo}

Como o modelo de referência é um modelo de linguagem autorregressivo, a métrica de
desempenho preditivo adotada é a \textit{perplexity}, padrão para essa classe de modelos.
Ela quantifica o grau de ``surpresa'' do modelo diante de cada \textit{token} do conjunto
de avaliação, constituindo uma medida da qualidade da modelagem probabilística da
linguagem \cite{Goodfellow2016}. Valores mais baixos de \textit{perplexity} indicam melhor
capacidade preditiva.

A adoção da \textit{perplexity}, em lugar de métricas de tarefas descendentes --- como a
acurácia em uma tarefa de classificação ---, decorre de ela ser uma métrica intrínseca:
pode ser medida diretamente sobre o corpus de avaliação, sem exigir o \textit{fine-tuning}
do modelo em uma tarefa específica, procedimento que tornaria a comparação entre
configurações dependente de um segundo treinamento. O procedimento de cálculo adotado é
descrito no Capítulo~\ref{chap:desenvolvimento}.

\subsection*{Custo computacional}

O custo computacional é avaliado por meio de quatro indicadores complementares:

\begin{itemize}
    \item \textbf{Número de parâmetros remanescentes:} contagem direta dos tensores de pesos
          não-zerados (no caso de poda não estruturada) ou existentes (no caso de poda
          estruturada) no modelo após a aplicação da técnica;

    \item \textbf{Operações de ponto flutuante (FLOPs):} contagem das operações matemáticas
          necessárias para uma inferência completa, obtida por análise estática do grafo
          computacional do modelo. Bibliotecas específicas --- como \texttt{thop}, em
          ambiente PyTorch --- percorrem cada camada do modelo e acumulam as operações
          correspondentes. A automação, porém, é parcial: a contagem cobre apenas os tipos
          de camada previstos no catálogo da ferramenta, de modo que camadas fora dele
          exigem o registro de regras de contagem próprias, sob pena de serem silenciosamente
          ignoradas. O tratamento efetivamente necessário no modelo adotado, bem como as
          omissões remanescentes, são detalhados no Capítulo~\ref{chap:desenvolvimento};

    \item \textbf{Tempo de inferência:} medido empiricamente, registrando o tempo decorrido
          durante o processamento de uma entrada representativa em condições controladas de
          \textit{hardware}. Múltiplas medições são agregadas por meio de média e
          desvio-padrão para mitigar a variância natural de execução;

    \item \textbf{Tamanho do modelo:} avaliado em duas dimensões --- espaço ocupado em disco
          após serialização do estado do modelo e consumo de memória durante a inferência,
          este último registrado por ferramentas de \textit{profiling} do PyTorch.
\end{itemize}

\subsection*{Impacto ambiental}

O impacto ambiental é mensurado por meio de dois indicadores:

\begin{itemize}
    \item \textbf{Consumo energético:} estimado por meio de ferramentas de monitoramento do
          uso de \textit{hardware} (CPU e GPU) durante a execução dos experimentos. A
          biblioteca CodeCarbon \cite{CodeCarbon2023} integra leituras periódicas de consumo
          elétrico dos componentes, registrando a energia total consumida (em kWh) durante
          o intervalo de inferência;

    \item \textbf{Emissões equivalentes de CO\textsubscript{2}:} calculadas a partir do
          consumo energético registrado, multiplicado por um fator de intensidade carbônica
          da matriz energética da região onde os experimentos são executados. A biblioteca
          CodeCarbon mantém base desses fatores e realiza essa conversão automaticamente.
          Cabe uma ressalva sobre o alcance desse indicador: como o fator de intensidade
          carbônica depende da localidade em que o código executa, ele não é controlado pelo
          experimento e pode variar entre execuções. As emissões devem, por isso, ser lidas
          como aproximações de ordem de grandeza, e a comparação entre configurações
          apoia-se no consumo energético, não nas emissões dele derivadas.
\end{itemize}

% Todas as métricas são calculadas tanto para o modelo de referência quanto para cada
% configuração de modelo podado, de modo que os resultados expressem ganhos ou perdas
% proporcionais à compressão introduzida, conforme detalhado a seguir.
}

% -------------------------------------------------------
% AJUSTE: reunião 13/05 — reescrita completa da Seção 3.5 (como medir e analisar energia)
% \section{Análise de Consumo Energético}\label{sec:metodologia_energia}
% % -------------------------------------------------------

% \adicionadomarcelo{Esta seção detalha o procedimento de coleta e análise das métricas de consumo energético,
% complementando as definições apresentadas na seção anterior.

% \subsection*{Mensuração do consumo energético absoluto}

% A medição do consumo energético é realizada por meio de ferramentas de \textit{software}
% que monitoram o uso de \textit{hardware} durante a execução dos experimentos. Em
% particular, a biblioteca CodeCarbon \cite{CodeCarbon2023} integra leituras periódicas de
% consumo elétrico de CPU e GPU, registrando a energia total consumida ao longo de cada
% sessão de inferência sobre o conjunto de teste. O resultado é uma medida de energia em
% kWh, acompanhada de uma estimativa de emissões equivalentes de
% CO\textsubscript{2}.

% Para mitigar a variância natural das medições de energia --- que tendem a ser mais ruidosas
% do que medições de desempenho preditivo ---, cada configuração de modelo é avaliada em
% múltiplas execuções com sementes aleatórias distintas. Os resultados são reportados como
% média e desvio-padrão.

% \subsection*{Cálculo da redução relativa}

% Para cada configuração de modelo podado, a redução relativa de consumo energético é
% calculada por comparação direta com o consumo do modelo de referência, conforme a
% expressão:

% \[
%     \text{Redução relativa} = \frac{E_{\text{baseline}} - E_{\text{podado}}}{E_{\text{baseline}}} \times 100\%
% \]

% \noindent em que $E_{\text{baseline}}$ e $E_{\text{podado}}$ representam, respectivamente,
% a energia consumida pelo modelo de referência e pelo modelo podado durante a inferência
% sobre o mesmo conjunto de teste. Todas as comparações são realizadas em iguais condições
% de \textit{hardware}, mesmo conjunto de teste e mesmo número de iterações de inferência,
% de modo a garantir a validade das comparações.

% \subsection*{Relação entre FLOPs, parâmetros e consumo energético}

% A redução teórica de FLOPs, calculada por análise estática do grafo computacional,
% constitui um indicador \textit{proxy} da redução de custo computacional, mas não se
% traduz linearmente em redução de energia. \textit{Hardware} moderno apresenta componentes
% de consumo basal --- como memória RAM, controladores e sistemas de refrigeração --- cujo
% consumo não escala proporcionalmente com o número de operações realizadas.

% A redução de parâmetros impacta o tamanho do modelo em memória, o que pode diminuir as
% transferências entre memória e processador --- fonte significativa de consumo energético
% em arquiteturas de \textit{hardware} contemporâneas. A poda estruturada tende a apresentar
% redução de energia proporcional à redução de FLOPs, pois o \textit{hardware} efetivamente
% executa menos operações. A poda não estruturada, em contrapartida, pode reduzir os FLOPs
% teóricos sem reduzir proporcionalmente o consumo energético real, uma vez que, em
% \textit{hardware} sem suporte nativo a esparsidade, multiplicações por zero ainda consomem
% ciclos de processamento.

% Essa análise sustenta a comparação central deste trabalho: validar empiricamente se
% diferentes paradigmas de poda traduzem-se ou não em ganhos energéticos reais, e em que
% proporção.

% \subsection*{Significado prático em larga escala}

% Modelos de linguagem de grande escala em produção servem milhões de requisições diárias.
% Nessa escala, mesmo reduções modestas no consumo energético por inferência podem resultar
% em economia significativa de energia e redução de emissões de carbono ao longo do tempo
% \cite{Strubell2019}. A análise conduzida neste trabalho, embora realizada em escala
% experimental, busca fornecer evidências quantitativas que permitam estimar o impacto
% potencial dessas técnicas quando aplicadas em cenários de produção.}
